\documentclass[12pt,a4paper]{article}
\usepackage[T2A]{fontenc}
\usepackage[utf8]{inputenc}
\usepackage[russian,english]{babel}
\usepackage{geometry}
\usepackage{graphicx}
\usepackage{booktabs}
\usepackage{longtable}
\usepackage{float}
\usepackage{array}
\usepackage{hyperref}
\usepackage{setspace}
\geometry{margin=1in}
\hypersetup{colorlinks=true,linkcolor=black,urlcolor=blue}
\setlength{\parindent}{1.25cm}
\setlength{\parskip}{0.35em}
\onehalfspacing

\begin{document}

\begin{titlepage}
\begin{center}
{\large Assignment №2}\\
{\large Data-Driven Decision Project (stage 2)}\\[2cm]
{\LARGE \textbf{Data Preparation and Market Intelligence}}\\[1cm]
{\Large \textbf{Staj.kz --- Платформа стажировок для студентов Казахстана}}\\[3cm]
\end{center}

\noindent \textbf{Имя и фамилия:} Азамат Ержанулы, Данияр Койшин\\
\textbf{Академическая группа:} CSE-2503M\\
\textbf{Дата:} 03.04.2026

\vfill

\begin{center}
Алматы, 2026
\end{center}
\end{titlepage}

\section*{Введение}
Во втором этапе проекта нужно было собрать и подготовить данные о конкурентах, а затем на их основе оценить, в какой сегмент рынка платформе Staj.kz лучше заходить. В первом задании выборка получилась слишком маленькой, поэтому в этой работе рынок был рассмотрен шире: не только стажировки, но и trainee-, graduate- и junior-позиции, а также вакансии без требования опыта.

Для анализа использовались два источника: hh.kz как крупнейшая коммерческая платформа вакансий и Enbek.kz как государственная платформа занятости. Данные собирались программно, затем объединялись в одну таблицу, очищались от дублей, проверялись на полноту и свежесть, а зарплатные значения дополнительно проверялись на выбросы методом 1.5xIQR. Все основные изменения были зафиксированы в журнале очистки.

\section{Задание 1. Сбор, очистка и объединение данных}

\subsection{Источники данных и логика отбора}
\begin{itemize}
\item Использован источник \texttt{hh.kz / api.hh.ru} как основной коммерческий рынок вакансий.
\item Использован источник \texttt{Enbek.kz} как государственная платформа занятости, включая youth- и practice-маршруты.
\item В выборку включались не только стажировки, но и более широкий early-career сегмент: internship, trainee, graduate, junior, no-experience, youth, practice.
\item География исследования --- весь Казахстан.
\end{itemize}

\subsection{Краткая характеристика итогового датасета}
\begin{itemize}
\item Сырые записи до очистки: \texttt{3881}
\item Итоговые записи после очистки: \texttt{2538}
\item Доля вакансий с указанной зарплатой: \texttt{75.49\%}
\item Доля вакансий, релевантных студентам и выпускникам: \texttt{87.04\%}
\item Окно свежести вакансий: \texttt{90} дней
\end{itemize}

\subsection{Сравнение источников}
\begin{longtable}{p{0.16\linewidth}p{0.16\linewidth}p{0.18\linewidth}p{0.38\linewidth}}
\toprule
Источник & Количество вакансий & Доля вакансий с зарплатой, \% & Характеристика\\
\midrule
\endfirsthead
\toprule
Источник & Количество вакансий & Доля вакансий с зарплатой, \% & Характеристика\\
\midrule
\endhead
hh.kz & 2347 & 76.35 & Коммерческая платформа с высокой шириной рынка и сильным поиском.\\
Enbek.kz & 191 & 64.92 & Государственная платформа с отдельными youth/practice-маршрутами.\\
\bottomrule
\end{longtable}

\subsection{Документирование очистки данных}
Очистка данных выполнялась в несколько этапов:
\begin{enumerate}
\item удаление дублей по полям источника и идентификатора вакансии;
\item дополнительная мягкая дедупликация по сочетанию названия вакансии, компании, города и даты публикации;
\item проверка полноты и исключение записей без критически важных полей;
\item нормализация городов, графиков работы, типов занятости и валют;
\item приведение зарплат к единому показателю \texttt{salary\_mid\_kzt};
\item проверка выбросов по методу 1.5xIQR и обнуление явно аномальных значений.
\end{enumerate}

Отдельно проверялись пропуски, различия в написании городов и категорий занятости, лишние пробелы и другие мелкие несоответствия, которые мешали бы корректному объединению данных из двух источников.

\begin{longtable}{p{0.20\linewidth}p{0.50\linewidth}p{0.16\linewidth}}
\toprule
Этап & Метрика & Значение\\
\midrule
\endfirsthead
\toprule
Этап & Метрика & Значение\\
\midrule
\endhead
Сырые данные & Количество записей & 3881.0\\
Дедупликация & После удаления дублей по источнику и идентификатору вакансии & 2574.0\\
Проверка свежести & Удалено устаревших или некорректно датированных записей & 0.0\\
Проверка полноты & После фильтра по обязательным полям & 2574.0\\
Дедупликация & Удалено мягких дублей & 36.0\\
Проверка выбросов & Зарплат обнулено как выбросы & 166.0\\
Проверка выбросов & Нижняя граница IQR, KZT & -337500.0\\
Проверка выбросов & Верхняя граница IQR, KZT & 1122500.0\\
Сегментация & Q25 для сегментации, KZT & 200000.0\\
Сегментация & Q50 для сегментации, KZT & 290000.0\\
Сегментация & Q75 для сегментации, KZT & 475000.0\\
Итог & Количество записей & 2538.0\\
\bottomrule
\end{longtable}

После объединения двух источников было получено 3881 сырых записей. После удаления дублей, нормализации и обработки выбросов итоговый датасет составил 2538 записей, то есть требование задания по объёму выполнено. Этого объёма достаточно для дальнейшего анализа рынка и сравнения конкурентов.

\section{Задание 2. Анализ конкурентной среды и рыночные выводы}

Анализ конкурентной среды проводился в двух частях. Первая часть опиралась на очищенный датасет вакансий и показывала, как устроен рынок по городам, ролям и зарплатам. Вторая часть была посвящена сравнению самих платформ как конкурентов Staj.kz.

\subsection{География рынка}
\begin{longtable}{p{0.45\linewidth}p{0.22\linewidth}}
\toprule
Город & Количество вакансий\\
\midrule
\endfirsthead
\toprule
Город & Количество вакансий\\
\midrule
\endhead
Алматы & 924\\
Астана & 586\\
Караганда & 140\\
Шымкент & 113\\
Костанай & 102\\
Усть-Каменогорск & 63\\
Павлодар & 60\\
Актау & 54\\
\bottomrule
\end{longtable}

Наибольшая концентрация вакансий наблюдается в Алматы и Астане. Это означает, что запуск Staj.kz логичнее начинать именно с этих городов: здесь выше и число вакансий, и вероятность быстрее собрать начальную базу работодателей и кандидатов.

\subsection{Структура спроса по категориям ролей}
\begin{longtable}{p{0.48\linewidth}p{0.22\linewidth}}
\toprule
Категория роли & Количество вакансий\\
\midrule
\endfirsthead
\toprule
Категория роли & Количество вакансий\\
\midrule
\endhead
Маркетинг и продажи & 919\\
Общий early-career сегмент & 767\\
Образование и поддержка & 396\\
Финансы и административные роли & 191\\
Операции и логистика & 162\\
IT и аналитика данных & 103\\
\bottomrule
\end{longtable}

По структуре рынка видно, что сильнее всего представлены категории \textit{маркетинг и продажи} и \textit{общий early-career сегмент}. Это важно для Staj.kz, потому что продукт на старте не должен ограничиваться только IT-стажировками: заметная часть спроса находится и в более массовых направлениях.

\subsection{Ценовые сегменты}
Ценовые сегменты определялись по квартилям распределения \texttt{salary\_mid\_kzt}. Это позволило отказаться от субъективного назначения границ и сегментировать рынок на основе фактических данных.

\begin{longtable}{p{0.40\linewidth}p{0.22\linewidth}}
\toprule
Ценовой сегмент & Количество вакансий\\
\midrule
\endfirsthead
\toprule
Ценовой сегмент & Количество вакансий\\
\midrule
\endhead
не указан & 622\\
низкий (low-priced) & 487\\
высокий (high-priced) & 480\\
средний (middle-priced) & 477\\
премиальный (luxury) & 472\\
\bottomrule
\end{longtable}

По итогам анализа наиболее перспективным для запуска Staj.kz оказался сегмент \texttt{средний (middle-priced)}. В нём достаточно вакансий, при этом он выглядит более реалистичным для коммерческого запуска, чем слишком узкий премиальный сегмент или слишком размытый низкий.

\subsection{Крупнейшие работодатели}
\begin{longtable}{p{0.68\linewidth}r}
\toprule
Работодатель & Количество вакансий\\
\midrule
\endfirsthead
\toprule
Работодатель & Количество вакансий\\
\midrule
\endhead
Работа-Подработка: Курьер, Кладовщик, Сборщик & 57\\
“QSR” ,ТМ BURGER KING & 50\\
TECHNODOM Operator (Технодом Оператор) & 45\\
Magnum E-commerce Kazakhstan & 38\\
ЦАН & 33\\
Сапаров & 32\\
Этажи & 31\\
Яндекс & 24\\
Andersen & 20\\
Августа & 18\\
\bottomrule
\end{longtable}

\subsection{Извлечённые навыки}
\begin{longtable}{p{0.68\linewidth}r}
\toprule
Навык & Частота\\
\midrule
\endfirsthead
\toprule
Навык & Частота\\
\midrule
\endhead
Ответственность; & 49\\
Аккуратность; & 39\\
Адаптивность; & 36\\
Дисциплинированность; & 32\\
Вежливость; & 22\\
Выносливость и усердие; & 17\\
Добропорядочность; & 16\\
Логическое мышление; & 16\\
Самостоятельность и ответственность; & 14\\
Компьютерная грамотность; & 14\\
\bottomrule
\end{longtable}

\subsection{Ключевые потребители конкурентов}
На основе структуры данных были выделены следующие типы потребителей конкурентных платформ:
\begin{itemize}
\item студент старших курсов, ищущий первую стажировку или позицию начального уровня;
\item выпускник без опыта, ориентированный на junior-позиции;
\item HR малого и среднего бизнеса, которому важна скорость закрытия вакансии и стоимость размещения;
\item HR крупной компании, которому важны фильтры, аналитика и качество откликов.
\end{itemize}

\subsection{Сравнение конкурентов по технологическому уровню}
\begin{longtable}{p{0.15\linewidth}p{0.16\linewidth}p{0.16\linewidth}p{0.16\linewidth}p{0.16\linewidth}p{0.19\linewidth}}
\toprule
Платформа & Фокус на студентах & Поиск и фильтры & Прозрачность зарплаты & Технологический уровень & Комментарий\\
\midrule
\endfirsthead
\toprule
Платформа & Фокус на студентах & Поиск и фильтры & Прозрачность зарплаты & Технологический уровень & Комментарий\\
\midrule
\endhead
hh.kz & Средний & Высокий & Средняя & Высокий & Сильная универсальная платформа, но не специализирована под студентов.\\
Enbek.kz & Средне-высокий & Средний & Высокая & Средний & Сильна за счёт государственных сценариев занятости и youth/practice-категорий.\\
Staj.kz & Очень высокий & Нишевой таргетинг & Очень высокая & Ориентация на мобильный формат & Может занять нишу сервиса, ориентированного именно на студентов и выпускников без опыта.\\
\bottomrule
\end{longtable}

\subsection{SWOT-анализ}
\textbf{Сильные стороны}
\begin{itemize}
\item итоговый датасет содержит 2538 записей, поэтому требование задания по объёму выполнено;
\item данные из двух источников дают более полную картину рынка, чем выборка только с одного сайта;
\item большая доля вакансий без опыта подтверждает, что идея Staj.kz опирается на реальный спрос.
\end{itemize}

\textbf{Слабые стороны}
\begin{itemize}
\item часть вакансий публикуется без зарплаты, поэтому студентам сложнее сравнивать предложения;
\item крупные платформы охватывают рынок широко, но не дают узкой настройки именно под студенческий сегмент;
\item не весь рынок состоит из буквальных стажировок, поэтому позиционирование Staj.kz нужно формулировать шире.
\end{itemize}

\textbf{Возможности}
\begin{itemize}
\item сделать платформу более удобной для студентов за счёт прозрачной зарплаты и понятной фильтрации;
\item сфокусироваться на среднем ценовом сегменте как на наиболее устойчивом для старта;
\item начинать запуск с Алматы и Астаны, где рынок уже достаточно плотный.
\end{itemize}

\textbf{Угрозы}
\begin{itemize}
\item универсальные платформы уже собирают большую часть трафика вакансий;
\item рынок зависит от сезонности найма и от того, что многие вакансии остаются очными;
\item нишевой платформе нужно доказать работодателям, что качество откликов у неё действительно выше.
\end{itemize}

\section{Дашборды и визуализации}
\begin{figure}[H]
\centering
\includegraphics[width=0.92\textwidth]{\detokenize{../figures/01_source_share.png}}
\caption{Распределение вакансий по источникам}
\end{figure}
\begin{figure}[H]
\centering
\includegraphics[width=0.92\textwidth]{\detokenize{../figures/02_top_cities.png}}
\caption{Города-лидеры по количеству вакансий}
\end{figure}
\begin{figure}[H]
\centering
\includegraphics[width=0.92\textwidth]{\detokenize{../figures/03_role_categories.png}}
\caption{Структура рынка по категориям ролей}
\end{figure}
\begin{figure}[H]
\centering
\includegraphics[width=0.92\textwidth]{\detokenize{../figures/04_price_segments.png}}
\caption{Распределение вакансий по ценовым сегментам}
\end{figure}
\begin{figure}[H]
\centering
\includegraphics[width=0.92\textwidth]{\detokenize{../figures/05_publication_timeline.png}}
\caption{Динамика публикации вакансий}
\end{figure}
\begin{figure}[H]
\centering
\includegraphics[width=0.92\textwidth]{\detokenize{../figures/06_salary_transparency.png}}
\caption{Прозрачность зарплат по источникам}
\end{figure}

\section*{Заключение}
В результате анализа можно сделать вывод, что для Staj.kz наиболее разумно начинать со среднего ценового сегмента. Он даёт достаточный объём вакансий и при этом выглядит более понятным для коммерческого запуска. Основная идея платформы должна строиться не на попытке конкурировать со всеми сразу, а на понятной специализации: студенты и выпускники без опыта, прозрачные зарплаты, удобный отклик и более точная фильтрация вакансий.

\end{document}
